\section{The kernel and smooth even continuation}\label{sec:kernel}

Put
\[
 \vartheta(x)=\sum_{n\in\mathbb Z}e^{-\pi n^2x},\qquad
 \Theta(t)=\vartheta(e^{2t}),\qquad H(t)=e^{t/2}\Theta(t).
\]
The theta transformation \(\vartheta(x)=x^{-1/2}\vartheta(1/x)\) gives
\[
 \Theta(-t)=e^t\Theta(t),\qquad H(-t)=H(t).
\]
The series and all its derivatives converge locally uniformly, so \(H\) is
real analytic and even. Define
\[
 \Phi(t)=\frac12\left(H''(t)-\frac14H(t)\right).
\]
This is twice the common de Bruijn kernel normalization; multiplication by a
positive constant changes none of the PF signs or strictness statements.
Termwise differentiation gives, for every real \(t\),
\begin{equation}\label{eq:kernel-series}
 \Phi(t)=\sum_{n\geq1}
 \left(4\pi^2n^4e^{9t/2}-6\pi n^2e^{5t/2}\right)
 e^{-\pi n^2e^{2t}}.
\end{equation}
Thus the positive-side formula often written with \(|t|\) is the restriction
of a real-analytic even function. In particular, smoothness at the origin is
not inferred from numerical coverage.

For \(t\geq0\), the coefficient of the \(n\)-th exponential can be factored as
\[
 2\pi n^2e^{5t/2}\bigl(2\pi n^2e^{2t}-3\bigr)>0,
\]
because \(2\pi>3\). Hence every term in \eqref{eq:kernel-series} is positive
and \(\Phi(t)>0\) on \([0,\infty)\). Evenness gives \(\Phi(t)>0\) on all of
\(\R\), so the following logarithm is globally defined.

Write
\[
 \ell=\log\Phi,\qquad s=\ell',\qquad q=-s'=-\ell''.
\]
For \(a<b\), define
\begin{equation}\label{eq:AM}
 A(a,b)=s(a)-s(b),\qquad
 M(a,b)=\frac{q(b)-q(a)}{A(a,b)}.
\end{equation}
When \(q>0\), \(s\) is strictly decreasing and \(A(a,b)>0\).
